\subsection{Assumptions} \label{assumptions}

\begin{figure*}[t]
\begin{pchstack}[boxed,center,space=1em]

\begin{pchstack}
  \procedure{$\main$ $\mathcal{G}^{\dl}_{\A}(\secp)$}{
  (\Gr, p , g) \randpick \GroupGen(\usecp) \\
  x \randpick \Zq;~X \gets g^{x} \\
  x' \randpick \A((\Gr, p , g), X) \\
	\pcif x' = x \\
	\pcind \pcreturn 1\\
	\pcreturn 0
  }
\end{pchstack}

\begin{pchstack}
  \vline
      \pchspace
  \procedure{$\main$ $\G^{\ell\hbox{-}\colorbox{gray}{$\mathsf{a}$}\omdl}_{\A}(\secp)$ }{
  (\Gr, p , g) \randpick \GroupGen(\usecp) \\
  Q \gets \emptyset \\
  q \gets 0 \\
  \pcfor i \in [0..\ell] \pcdo \\
  \pcind x_i \randpick \Zq;~X_i \gets g^{x_i} \\
  \pcind Q[X_i] = x_i \\
  \vec{x} \gets (x_0, \ldots, x_\ell) \\
  \vec{X} \gets (X_0, X_1, \dots, X_{\ell}) \\
  \vec{x'} \gets \A^{\BigO^{\dl}}((\Gr, p, g), \vec{X})  \\
  	\pcif \vec{x'} = \vec{x}~\wedge~q < \ell \\
	\pcind \pcreturn 1\\
	\pcreturn 0
  }
  \pchspace
  \procedure{$\BigO^{\dl}(X, \colorbox{gray}{$\alpha, \{\beta_i\}_{i=0}^{\ell}$} )$\  }{
    \pccomment{\colorbox{gray}{$X = g^{\alpha} \prod_{i=0}^{\ell} X_i^{\beta_i}$}} \\
  q \gets q + 1 \\
  %\colorbox{light-gray}{$\pcreturn \bot\ \pcif Q[X] = \bot$} \\
  \colorbox{gray}{$x \gets \alpha + \sum_{i=0}^{\ell} x_i \beta_i$} \\
  \colorbox{gray}{$\pcreturn x$} \\
  x \gets \mathsf{dlog}(X) \\
  \pcreturn x
  }
\end{pchstack}

\end{pchstack}
    \caption{The Discrete Logarithm (DL), One-More Discrete Logarithm (OMDL), and Algebraic One-More Discrete Logarithm (AOMDL) games.
    $\Gr$ is a cyclic group with prime order $q$ and generator $g$.
    $\mathsf{dlog}$ is an algorithm that finds the discrete logarithm relation of the input $X$ and $g$.
    The differences between the OMDL and AOMDL games are highlighted in gray;
    the key distinction is that in the AOMDL game,
    the adversary can only query for linear combinations of its challenge group elements;
    in the OMDL game, the environment must return the discrete logarithm of an
    arbitrary group element.}
    \label{fig:dl-omdl}
\end{figure*}

\paragraph{Discrete Logarithm (DL) Assumption.}
I show the DL assumption along with the DL assumption in Figure~\ref{fig:dl-omdl},
but give an intuition here.
To begin,
a PPT adversary is given a value $X \in \GG$ such that $X=g^x$ without knowing the value of $x$.
The challenge to the adversary is to output $x$.

We say that the DL assumption holds if the advantage to a PPT adversary in
successfully producing such an $x$ is a negligible function of $\lambda$.

\paragraph{Computational Diffie-Hellman Assumption (CDH).}
The Computational Diffie-Hellman Assumption (CDH) can be described as follows.
A PPT adversary is given the following values:

\[
  (g, g^a, g^b)
\]

The challenge to the adversary is then to compute $g^{ab}$.
The CDH assumption holds if the advantage to any PPT adversary in successfully
doing so is a negligible function of $\lambda$.

\paragraph{One-More Discrete Logarithm Assumption (OMDL).}
The One-More Discrete Logarithm (OMDL) assumption~\cite{BellareNPS03,BauerFP21} extends
the DL assumption, which is non-interactive, into an interactive assumption.
I show the OMDL assumption along with the DL assumption in Figure~\ref{fig:dl-omdl}.

In particular,
a PPT adversary is given $\ell+1$ discrete logarithm challenges $(X_0, \ldots, X_\ell) \in \GG^{\ell+1}$
with respect to a generator $g \in \GG$.
It is also given access to a discrete logarithm oracle,
that when queried on an arbitrary group element $X' \in \GG$,
returns the discrete logarithm of that value.
The adversary is said to win the OMDL game if it makes $\ell$ or fewer queries
to the discrete logarithm oracle,
and produces correct solutions $(x_0, \ldots, x_\ell)$ to its $\ell+1$ challenges such that $X_i = g^{x_i}$,
$i \in \{ 0, \ldots, \ell\}$.

The OMDL assumption holds if the advantage to any PPT adversary in winning the
OMDL game is a negligible function of $\lambda$.

\paragraph{Algebraic One-More Discrete Logarithm Assumption (AOMDL).}
The algebraic one-more discrete logarithm assumption (AOMDL) was formalized by Jonas, Ruffing, and Seurin~\cite{NickRS21}.
We highlight differences with the OMDL assumption in Figure~\ref{fig:dl-omdl}.
Note that while the OMDL assumption is not falsifiable
(because the adversary is allowed to query for the discrete logarithm of any arbitrary group element).
the AOMDL assumption is falsifiable.
This is because the input to the discrete logarithm oracle is a linear combination of all challenges issued to the adversary,
so it is guaranteed that the environment runs in polynomial time with respect to a PPT adversary.
Hence, the AOMDL assumption is a strictly weaker assumption than standard OMDL.

An adversary in the AOMDL game is initialized in exactly the same manner as an OMDL adversary,
with the same winning condition.
The only distinction between the games is the inputs to the discrete logarithm oracle.
In the AOMDL setting, the adversary provides a vector $(\alpha, \beta_0 \ldots, \beta_\ell)$ of coefficients to the discrete logarithm oracle.
The oracle then responds with the integer linear combination of these coefficients with respect to the
discrete logarithm of the set of challenges $(X_0, \ldots, X_\ell)$.
Informally, an AOMDL adversary can be used to win the OMDL game,
therefore implying that if OMDL is hard, then AOMDL is hard.

The AOMDL assumption holds if the advantage to any PPT adversary in winning the
AOMDL game is a negligible function of $\lambda$.

\paragraph{Algebraic Group Model (AGM).}
Introduced by Fuchsbauer, Kiltz, and Loss~\cite{FuchsbauerKL18},
the Algebraic Group Model (AGM) gives an idealized representation of a group,
by requiring the adversary to act algebraically.
In other words,
when an adversary outputs a group element $X \in \GG$ to the environment,
it additionally is required to output its representation with respect to all
other group elements the adversary has seen thus far.
In other words, the adversary must output a coefficient vector $\textbf{v} = (\alpha, b_1, b_2, \ldots)$,
such that $X = g^\alpha A_1^{b_1} A_2^{b_2} \ldots$.


\paragraph{Knowledge of Exponent Assumption (KoE).}

The Knowledge of Exponent (KoE) assumption~\cite{Damgard91,BellareP04} is an unfalsifiable assumption.
In short,
the assumption says that for an adversary that is given input $(g, h=g^s, q)$,
if the adversary outputs a value $(C, Y) \in \GG^2$ such that $C^s = Y$,
then there exists an extractor algorithm that given $(g, h=g^s, q)$,
%the adversary's state and $(C, Y) \in \GG^2$,
the extractor can output $x$ such that $C = g^x$ (and, consequently $Y = C^x$).
Intuitively, KoE assumes that in order to output such a $(C, Y)$,
the adversary must know $x$.


